% ============================================================
%  Section 6 — RAG Engine
% ============================================================
\section{RAG Engine}
\label{sec:rag}

\subsection{Retrieval Design}

The Coder uses retrieval to ground account assignment in a company's own
accounting rules. The local knowledge base contains a chart of accounts and
topic-level Markdown policy sources per company. Inference embeds a query built
from supplier, document type, and description; Qdrant returns up to five chunks
filtered by \texttt{company\_id}. Retrieved text is passed alongside the
company's versioned Coder skill. The embedding model is
\texttt{paraphrase-multilingual-mpnet-base-v2} (768 dimensions), and the vector
store runs locally in embedded mode.

The implementation supports two chunkers. The default fixed-window strategy
uses 512 tokens with a 64-token overlap. The heading-based strategy groups
Markdown sections up to 1024 tokens. These choices are recorded for comparison;
the present corpus is too small to establish a winning strategy.

\subsection{What the Retrieval Evaluation Establishes}

Twenty hand-written queries per company were evaluated against expected policy
source filenames. With fixed chunks, source-level Recall@5 is 1.00 for both
companies; MRR@5 is 0.9167 for Company A and 0.8875 for Company B. With heading
based chunks, Recall@5 remains 1.00; MRR@5 is 0.9167 and 0.8500, respectively.
The fixed corpus contains 9 Company A chunks and 11 Company B chunks; the
heading corpus contains 9 and 10.

These figures should not be read as rule-level retrieval accuracy. The corpus
is still small relative to top-k=5, and the label only checks whether the
expected policy source appears; it does not identify the correct rule within
that source. Recall@5 is therefore saturated, and this small comparison does
not support a chunking winner. The evaluation verifies that company-filtered
retrieval returns the expected source, but it does not show that the best rule
ranks first or that retrieval improves accounting decisions.

The coding ablation in \cref{tab:baseline_main} is separate evidence: on the
reported selection runs, Coder score with retrieval was 0.30 for Company A and
0.21 for Company B, versus 0.10 and 0.00 without retrieval. These are low
absolute scores on synthetic data and do not isolate retrieval quality from
model and prompt errors. A stronger next evaluation should split the knowledge
base into individual rule chunks, label relevant rule IDs for each query, and
report Recall@k, MRR, and per-query examples on a larger rule set.
